\documentclass[12pt,letterpaper]{article} 
\usepackage[T1]{fontenc}
\usepackage{microtype}
\usepackage{hyperref}

%--- Page layout
\usepackage[margin=0.6in]{geometry}
\pagestyle{empty}                  % no page number on a one-pager
\setlength{\parindent}{0pt}        % no paragraph indents
\setlength{\parskip}{4pt}          % small gap between paragraphs instead

%--- Colours 
\usepackage{xcolor}
\definecolor{accent}{RGB}{28,55,94}    % navy: name, section labels, bullets
\definecolor{muted}{RGB}{105,112,122}  % grey: secondary text

%--- Lists
\usepackage{enumitem}
\setlist[itemize]{leftmargin=1.1em, itemsep=2pt, topsep=2pt, parsep=0pt, label={\color{accent}\textbullet}}  % compact lists, navy bullets

%--- Section labels
\newcommand{\sect}[2][]{
\par\vspace{8pt}
{\color{accent}\small\bfseries\MakeUppercase{#2}}\hfill{\color{muted}\small#1}\par
\vspace{2pt}\nointerlineskip\vspace{2pt}
{\color{accent!35}\rule{\linewidth}{0.5pt}}
\par\vspace{-1pt}
}

% PDF metadata
\usepackage[hidelinks]{hyperref}
%\hypersetup{pdftitle={Mill Street Labs}, pdfauthor={Mill Street Labs}}

\begin{document}

{\fontsize{26}{28}\selectfont\bfseries\color{accent} Mill Street Labs}\par
\vspace{2pt}{\color{accent}\rule{\linewidth}{1.2pt}}\par
\vspace{4pt}

\textbf{We envision frictionless intelligence transfer and accumulation among people and agents.} We are building long-running agent populations that discover, share, and build on useful knowledge.

\sect[Team of 4, Harvard Class of '27]{Team}
  Diego Luca Gonzalez Gauss (Math 55 TA, Susquehanna, Kempner research), Kevin Matthew Liu (Jump Trading, ex-Optiver, Kempner research), Enrico Yao-Bate (Math 55 Head TA, statistics research), \\ Haozhe Yang (Math 55 TA, Jane Street, IMO Medalist, Putnam Honorable Mentions)

\sect{Thesis}
\textbf{Long-running agent populations can improve collectively through adaptive social dynamics and intelligence transfer.}

Human innovation is an emergent property of long-running social dynamics; new ideas build on knowledge accumulated over generations through teaching, cooperation, and competition between groups.

\begin{itemize}
  \item \textbf{Adaptive social dynamics:} have groups compete for finite resources, so that helping one's group is rewarded even if doing so costs an agent resources. Let agents revise how they cooperate, enter and leave swarms, and migrate between groups.
  \item \textbf{Intelligence transfer:} credit agents for the improvement of those they teach; use teaching to combine discoveries from specialized lineages; accrue an interpretable corpus of verified communal knowledge with easy deposit and retrieval
\end{itemize}

\sect{Research} % I don't think research is the right title for this but edge also maybe isn't; agree but I think research sounds nicer lol
\textbf{Our focus is swarm improvement itself.}

Agent swarms have spontaneously demonstrated social dynamics, but their social development is limited by short-horizon objectives, siloed reasoning, and fixed search space.  We let populations run continually while perturbing organization and incentives, evaluating the impact on what agents discover and retain, to identify what sustains collective improvement.

\sect{Sandboxes}

\textbf{Sandbox 1}: strategy development in skillful games

Skillful games (e.g. predator-prey models, \href{https://arxiv.org/pdf/2609.19124}{Flag Game}, Werewolf, Hanabi) provide clear credit assignment, counterfactuals, and metrics for agent interaction.

\begin{itemize}
  \item developing custom policies, corpus, gestation, and fine-tuned personalities in SmolLM2 and GPT-4o
  \item accelerating SmolLM2 discovery fivefold in predator-prey environment over equivalent model calls
  \item saturating multiple skillful games with $\leq$ 1/10th the compute of baseline models
\end{itemize}

\textbf{Sandbox 2}: autonomous research in financial markets

Markets let swarms explore, build large ideas, innately reward ingenuity, and adapt to dynamic multivariate settings. We have domain knowledge to evaluate valuable problems and candidate strategies.

\begin{itemize}
  \item 3--4\,ms roundtrip execution with the Kalshi matching engine
  \item deployable positive expectancy strategies for market making
\end{itemize}

\end{document}
